1. Powers of i (cycle of 4)

$$ i^1 = i, \quad i^2 = -1, \quad i^3 = -i, \quad i^4 = 1 $$

To find $i^n$: compute r = n mod 4.

$$ \begin{array}{c|c}
r & i^r \\
0 & 1 \\
1 & i \\
2 & -1 \\
3 & -i
\end{array} $$

Example: 2026 = 4(506) + 2, so r = 2

$$ i^{2026} = i^2 = -1 $$

2. Sum of 4 consecutive powers of i equals zero

$$ i^k + i^{k+1} + i^{k+2} + i^{k+3} = 0 $$

3. Complex arithmetic

$$ (a + bi) + (c + di) = (a+c) + (b+d)i $$

$$ (a + bi)(c + di) = (ac - bd) + (ad + bc)i $$

$$ \frac{a + bi}{c + di} = \frac{(a+bi)(c-di)}{c^2 + d^2} = \frac{ac+bd}{c^2+d^2} + \frac{bc-ad}{c^2+d^2} \, i $$

$$ |a + bi| = \sqrt{a^2 + b^2} $$

4. Conjugate Root Theorem: If $a + bi$ is a zero of a real-coefficient polynomial, then $a - bi$ is also a zero.

5. Example: zeros 2, 3i, -3i

$$ P(x) = (x-2)(x-3i)(x+3i) = (x-2)(x^2+9) = x^3 - 2x^2 + 9x - 18 $$

So a = -2, b = 9, c = -18, and a + b + c = -11.

6. DeMoivre's Theorem

$$ \left( r(\cos\theta + i\sin\theta) \right)^n = r^n \left( \cos(n\theta) + i\sin(n\theta) \right) $$
